\section{Noninteger weights and the joint geometry of resonance}
\label{hqj:sec:rho}

For a complex parameter $\rho$, use $x^\rho=e^{\rho\log x}$ on the
positive spectrum. Let $\mathcal D(s,\rho)$ be the common-scale
linear-minus-quadratic difference with that weight. The meromorphic
continuation is most transparently written as
\begin{equation}\label{hqj:eq:rho-master}
\mathcal D(s,\rho)=\sum_{j\ge1}(-1)^jb^{2j}
C_j(s)\zeta(2s+2j-\rho,a),\qquad b<a.
\end{equation}
The same finite head as before globalizes this expression. In
particular, normal convergence is locally uniform in both $s$ and
$\rho$ away from the explicit polar divisors. The relevant positive
odd resonances are $\rho_0=2J-1$, $J\ge1$.

\begin{theorem}[Joint meromorphic normal form and finite-part conversion]
\label{hqj:thm:rho}
Near $(s,\rho)=(0,2J-1)$, with $\eta=\rho-(2J-1)$,
\begin{equation}\label{hqj:eq:rho-local}
\mathcal D(s,\rho)=
\frac{(-1)^Jb^{2J}C_J(s)}{2s-\eta}+G_J(s,\eta),
\end{equation}
where $G_J$ is jointly holomorphic and $G_J(s,\eta)=O(s^2)$.
For a nonresonant fixed $\rho$, define
$\Delta_r(\rho)=\partial_s^r\mathcal D(0,\rho)$.
Then $\Delta_1(\rho)=0$. At a resonant weight the fixed-$\rho$ derivative
instead satisfies
\begin{equation}\label{hqj:eq:rho-fp}
\boxed{\Delta_r(2J-1)
=\FP_{\rho=2J-1}\Delta_r(\rho)
+\frac{r!}{2}(-1)^Jb^{2J}c_{J,r+1}.}
\end{equation}
Here the finite part uses $\eta=\rho-(2J-1)$, and the off-resonance jets
are continued meromorphically in that variable. The extra finite term is
zero for $r\ge2J$.

Along a non-tangent ray $\eta=\kappa s$, $\kappa\ne2$, the first
coefficient is
\begin{equation}\label{hqj:eq:rho-ray}
\left.\frac{\dd}{\dd s}\mathcal D(s,2J-1+\kappa s)\right|_{s=0}
=\frac{(-1)^Jb^{2J}c_{J,2}}{2-\kappa}.
\end{equation}
The ray $\kappa=2$ lies on the polar divisor and requires subtraction
before restriction.
\end{theorem}
\begin{proof}
In \eqref{hqj:eq:rho-master}, exactly the $j=J$ term has a nearby
Hurwitz pole. Subtract $1/(2s-\eta)$ from that Hurwitz factor.
The remainder is jointly holomorphic. All other factors are locally
holomorphic, and normal convergence on a finite-head tail permits the
recombination. Every term, including the finite head, starts at $s^2$,
so $G_J=O(s^2)$. This proves \eqref{hqj:eq:rho-local} globally in the
positive mass, after the head procedure.

At $\eta\ne0$, the coefficient of $s^r$ in
$C_J(s)/(2s-\eta)$ is
\[
-\sum_{k=2}^r\frac{2^{r-k}c_{J,k}}{\eta^{r-k+1}}.
\]
It has no constant Laurent coefficient in $\eta$. At $\eta=0$, however,
the same meromorphic germ restricted before differentiation is
$C_J(s)/(2s)$, whose coefficient of $s^r$ is $c_{J,r+1}/2$.
The holomorphic remainder gives the same constant term in both orders
of operations. Multiplication by $r!(-1)^Jb^{2J}$ yields
\eqref{hqj:eq:rho-fp}. Since $\deg C_J=2J$, the last assertion about
vanishing of the local conversion follows. Restricting the normal form
to $\eta=\kappa s$ proves \eqref{hqj:eq:rho-ray}.
\end{proof}

\begin{remark}[Why a naive interpolation fails]
For nonresonant weights the first defect is zero, whereas at positive
odd weights it equals the nonzero local coefficient in
\eqref{hqj:eq:first-defect}. This is not a contradiction to analytic
continuation. Differentiation on different slices of a jointly
\emph{meromorphic} germ does not give a single jointly holomorphic
family across its intersecting divisor. At $\rho=-1$, the individual
spectra can themselves have poles at $s=0$ from their leading term;
that leading term cancels in $\mathcal D$. At that exceptional weight
we discuss derivatives of the common difference, not derivatives of
undefined individual values.
\end{remark}

The theorem resolves the resonant correction requested for $x^\rho$.
It also identifies exactly which local term is lost if the weight finite
part is taken before the spectral derivative. It does not assert
holomorphicity on the polar line or exchange two singular limits
without a subtraction convention.
